\section{Methodology and Data Collection}
\label{sec:methodology}

\subsection{Prompting and Code Generation}
\TODO{Define the prompt template; it differs between languages only in the target
language and the required I/O interface. Specify: zero-shot, no energy-related hints
(baseline), temperature, number of samples per task/LLM/language ($k$), handling of
non-compiling outputs (no manual repair, or a fixed number of automated retries).}

\subsection{Functional Correctness}
\TODO{Every generated program is executed against the task's test suite. Only programs
passing all tests enter the energy analysis. Correctness rates are reported separately
per LLM and language (\RQ{1}, \RQ{2}).}

\subsection{Measurement Procedure}
\TODO{Describe per run: warm-up, number of repetitions ($n$), randomised execution
order, cool-down / idle intervals, idle-power baseline, JVM warm-up policy, process
isolation, minimal background services. Record ambient temperature if possible.}

\subsection{Statistical Analysis}
\TODO{Planned analysis (to be finalised before data collection):
(i) descriptive statistics (median, IQR) per cell;
(ii) linear mixed-effects model on $\log(\text{energy})$ with fixed effects LLM,
language, platform and all interactions, and random intercepts for task (and generated
solution);
(iii) Type-III tests / likelihood-ratio tests for interactions, effect sizes, post-hoc
comparisons with multiplicity correction;
(iv) rank stability across languages and platforms via Kendall's $\tau$ / Friedman test
(\RQ{4});
(v) Spearman correlation between energy, time, and memory (\RQ{5}).
Check model assumptions; fall back to non-parametric or aligned-rank-transform methods
if violated.}
